\documentclass{article}

% --- Milestone 7 test: newtheorem definitions ---
\newtheorem{theorem}{Theorem}
\newtheorem{definition}{Definition}
\newtheorem{lemma}[theorem]{Lemma}
\newtheorem*{remark}{Remark}

% --- Milestone 7 test: newcommand with optional arg ---
\newcommand{\greeting}[1][World]{Hello, #1!}
\newcommand{\pair}[2][default]{(#1, #2)}

\begin{document}

\section{Introduction}
\label{sec:intro}
This is the intro section.

\section{Theorems}
\label{sec:thms}

\begin{theorem}
\label{thm:first}
Every even number greater than 2 is the sum of two primes.
\end{theorem}

\begin{lemma}
\label{lem:helper}
A useful helper lemma.
\end{lemma}

\begin{definition}
\label{def:prime}
A prime number is a natural number greater than 1.
\end{definition}

\begin{theorem}
\label{thm:second}
There are infinitely many primes.
\end{theorem}

\begin{remark}
This is an unnumbered remark.
\end{remark}

\section{Tables}
\label{sec:tables}

\begin{table}
\caption{A booktabs table}
\label{tab:booktabs}
\begin{tabular}{lcc}
\toprule
Name & Score & Grade \\
\midrule
Alice & 95 & A \\
Bob & 87 & B \\
\bottomrule
\end{tabular}
\end{table}

\section{Cross-references}
See \autoref{sec:intro} for the introduction.
See \autoref{thm:first} for the first result.
See \autoref{def:prime} for the definition.
See \autoref{tab:booktabs} for the table.
See \nameref{sec:intro} for details.

\section{Macros}
\greeting
\greeting[Alice]
\pair{second}
\pair[first]{second}

\end{document}
